% ═══════════════════════════════════════════════════════════════
% LERNBLATT: parecer y conocer (DS 1)
% Template: Best-Practice-Design v2.0
% ═══════════════════════════════════════════════════════════════
% DESIGN-PRINZIPIEN:
% - Sans-Serif (Helvetica): Fließtext
% - Serif (Palatino): Spanische Vokabeln (visuelle Distinktion)
% - Farbe NUR didaktisch begründet
% - Minimale Rahmen (kein "visual noise")
% ═══════════════════════════════════════════════════════════════

\documentclass[11pt, a4paper]{article}

\usepackage[utf8]{inputenc}
\usepackage[T1]{fontenc}
\usepackage[ngerman]{babel}

% --- SCHRIFTEN ---
\usepackage{helvet}
\renewcommand{\familydefault}{\sfdefault}
\usepackage{palatino}
\newcommand{\seriftext}[1]{{\fontfamily{ppl}\selectfont #1}}

\usepackage{microtype}
\usepackage[margin=1.8cm, top=2cm, bottom=1.5cm]{geometry}
\usepackage{enumitem}
\usepackage{tabularx}
\usepackage{booktabs}
\usepackage{array}
\usepackage{multicol}
\usepackage{tikz}
\usepackage{amssymb}

\usepackage{fancyhdr}
\usepackage{lastpage}

% ═══════════════════════════════════════════════════════════════
% VARIABLEN
% ═══════════════════════════════════════════════════════════════

\newcommand{\lektion}{1}
\newcommand{\thema}{parecer y conocer}
\newcommand{\untertitel}{Meinungen äußern und Personen kennen}
\newcommand{\klasse}{9}
\newcommand{\lernziele}{parecer wie gustar · conocer konjugieren · conocer + a bei Personen}

% ═══════════════════════════════════════════════════════════════
% FARBEN
% ═══════════════════════════════════════════════════════════════

\usepackage{xcolor}

\definecolor{hellgrau}{HTML}{F5F5F5}
\definecolor{rahmen}{HTML}{CCCCCC}
\definecolor{akzent}{HTML}{666666}
\definecolor{stamm}{HTML}{C62828}

% ═══════════════════════════════════════════════════════════════
% BOXEN (minimalistisch)
% ═══════════════════════════════════════════════════════════════

\usepackage{tcolorbox}
\tcbuselibrary{skins}

\newtcolorbox{lernzielbox}{
    colback=hellgrau, colframe=hellgrau, boxrule=0pt, arc=0pt,
    left=10pt, right=10pt, top=8pt, bottom=8pt,
    before skip=6pt, after skip=8pt
}

\newtcolorbox{grammatikbox}[1][]{
    colback=white, colframe=white, boxrule=0pt, arc=0pt,
    left=10pt, right=6pt, top=6pt, bottom=6pt,
    borderline west={3pt}{0pt}{akzent},
    before skip=4pt, after skip=4pt
}

\newtcolorbox{merkbox}[1][]{
    colback=hellgrau, colframe=hellgrau, boxrule=0pt, arc=0pt,
    left=10pt, right=10pt, top=8pt, bottom=8pt,
    borderline north={2pt}{0pt}{black},
    before skip=6pt, after skip=6pt
}

\newtcolorbox{vokabelbox}[1][]{
    colback=white, colframe=rahmen, boxrule=0.5pt, arc=0pt,
    left=8pt, right=8pt, top=6pt, bottom=6pt,
    before skip=4pt, after skip=4pt
}

\newtcolorbox{tippbox}{
    colback=white, colframe=white, boxrule=0pt, arc=0pt,
    left=8pt, right=4pt, top=3pt, bottom=3pt,
    borderline west={2pt}{0pt}{akzent}
}

\newtcolorbox{zusatzbox}[1][]{
    colback=white, colframe=rahmen, boxrule=0.5pt, arc=0pt,
    left=8pt, right=8pt, top=6pt, bottom=6pt,
    borderline north={1.5pt}{0pt}{akzent}
}

% ═══════════════════════════════════════════════════════════════
% BEFEHLE
% ═══════════════════════════════════════════════════════════════

\newcommand{\es}[1]{\seriftext{\textbf{#1}}}
\newcommand{\de}[1]{\textit{\textcolor{akzent}{#1}}}
\newcommand{\gram}[1]{\textsc{#1}}
\newcommand{\abmark}[1]{{\footnotesize\textcolor{akzent}{$\rightarrow$ AB #1}}}
\newcommand{\abschnitt}[2]{\noindent\textbf{\large #1. #2}}
\newcommand{\stw}[1]{\textcolor{stamm}{#1}}

\setlength{\parindent}{0pt}
\setlength{\parskip}{5pt}

% ═══════════════════════════════════════════════════════════════
% KOPF- UND FUSSZEILE
% ═══════════════════════════════════════════════════════════════

\pagestyle{fancy}
\fancyhf{}
\renewcommand{\headrulewidth}{0.5pt}
\renewcommand{\footrulewidth}{0pt}
\lhead{\footnotesize\textsc{Spanisch} · Klasse \klasse}
\rhead{\footnotesize\thema}
\cfoot{\footnotesize\thepage\,/\,\pageref{LastPage}}

% ═══════════════════════════════════════════════════════════════
% DOKUMENT
% ═══════════════════════════════════════════════════════════════

\begin{document}

% --- TITEL ---
\begin{center}
    {\LARGE\bfseries \thema}\\[3pt]
    {\small \untertitel}
\end{center}

\vspace{4pt}

% --- EINLEITUNG ---
Stell dir vor, dein spanischer Austauschschüler fragt dich:
\es{¿Qué te parece mi ciudad?} -- \de{Was hältst du von meiner Stadt?}
Oder du möchtest wissen, wen er kennt: \es{¿Conoces a mi hermano?}

\vspace{2pt}

% --- LERNZIELE ---
\begin{lernzielbox}
\textbf{Das lernst du:}\quad \lernziele
\end{lernzielbox}

% ═══════════════════════════════════════════════════════════════
% ZWEI SPALTEN
% ═══════════════════════════════════════════════════════════════

\begin{multicols}{2}

% ---------------------------------------------------------------
% ABSCHNITT 1: Vokabeln
% ---------------------------------------------------------------
\abschnitt{1}{Vocabulario} \de{-- Wortschatz}

\vspace{4pt}

\begin{tabular}{@{}ll@{}}
\es{parecer} & \de{scheinen, vorkommen} \\[0.4em]
\es{conocer} & \de{kennen, kennenlernen} \\[0.4em]
\midrule
\es{engreído/-a} & \de{eingebildet} \\[0.4em]
\es{simpático/-a} & \de{sympathisch} \\[0.4em]
\es{tranquilo/-a} & \de{ruhig} \\[0.4em]
\es{raro/-a} & \de{seltsam, komisch} \\[0.4em]
\es{guapo/-a} & \de{gutaussehend} \\[0.4em]
\midrule
\es{todavía no} & \de{noch nicht} \\[0.4em]
\es{bastante} & \de{ziemlich} \\
\end{tabular}

\vspace{4pt}

\abmark{Aufg. 1--2}

\vspace{8pt}

% ---------------------------------------------------------------
% ABSCHNITT 2: parecer
% ---------------------------------------------------------------
\abschnitt{2}{Gramática: parecer} \de{-- scheinen}

\vspace{4pt}

\begin{grammatikbox}
\textbf{\gram{parecer (wie gustar!)}}

\vspace{3pt}
{\small
\begin{tabular}{@{}ll@{}}
\es{me parece(n)} & \de{ich finde} \\[0.3em]
\es{te parece(n)} & \de{du findest} \\[0.3em]
\es{le parece(n)} & \de{er/sie findet} \\[0.3em]
\es{nos parece(n)} & \de{wir finden} \\[0.3em]
\es{os parece(n)} & \de{ihr findet} \\[0.3em]
\es{les parece(n)} & \de{sie finden} \\
\end{tabular}}
\end{grammatikbox}

\begin{tippbox}
\footnotesize\textit{Merke:} \es{parece} bei Singular, \es{parecen} bei Plural!
\end{tippbox}

\abmark{Aufg. 1}

\columnbreak

% ---------------------------------------------------------------
% ABSCHNITT 3: conocer
% ---------------------------------------------------------------
\abschnitt{3}{Gramática: conocer} \de{-- kennen}

\vspace{4pt}

\begin{grammatikbox}
\textbf{\gram{Konjugation}}

\vspace{3pt}
{\small
\begin{tabular}{@{}ll@{}}
yo & \es{cono\stw{zc}o} \\[0.3em]
tú & \es{conoces} \\[0.3em]
él/ella & \es{conoce} \\[0.3em]
nosotros & \es{conocemos} \\[0.3em]
vosotros & \es{conocéis} \\[0.3em]
ellos & \es{conocen} \\
\end{tabular}}

\vspace{4pt}
\textbf{Achtung:} yo = \es{cono\stw{zc}o} (c → zc)
\end{grammatikbox}

\vspace{4pt}

\begin{merkbox}
\textbf{conocer + a bei Personen!}

\vspace{4pt}
{\small
\textbf{Person} → mit \es{a}:\\
\es{¿Conoces \textbf{a} María?}

\vspace{3pt}
\textbf{Sache} → ohne a:\\
\es{¿Conoces la ciudad?}}
\end{merkbox}

\abmark{Aufg. 2--3}

\vspace{8pt}

% ---------------------------------------------------------------
% ABSCHNITT 4: Redemittel
% ---------------------------------------------------------------
\abschnitt{4}{Comunicación} \de{-- Redemittel}

\vspace{4pt}

\textbf{Meinung äußern:}\\
{\small
\es{¿Qué te parece...?} \de{Was meinst du zu...?}\\
\es{Me parece bien/mal.} \de{Ich finde es gut/schlecht.}\\
\es{Me parece interesante.} \de{Ich finde es interessant.}}

\abmark{Aufg. 4}

\end{multicols}

% ═══════════════════════════════════════════════════════════════
% SEITENUMBRUCH
% ═══════════════════════════════════════════════════════════════
\newpage

% ---------------------------------------------------------------
% ABSCHNITT 5: Beispieldialog
% ---------------------------------------------------------------
\abschnitt{5}{Diálogo modelo} \de{-- Beispieldialog}

\vspace{4pt}

\begin{merkbox}
\textbf{Über Personen sprechen}

\vspace{4pt}
{\small
\begin{tabular}{@{}p{7cm}p{6.5cm}@{}}
\textbf{A:} \es{¿Conoces a mi hermano Gonzalo?} & \de{Kennst du meinen Bruder Gonzalo?} \\[3pt]
\textbf{B:} \es{No, todavía no lo conozco. ¿Cómo es?} & \de{Nein, ich kenne ihn noch nicht. Wie ist er?} \\[3pt]
\textbf{A:} \es{Bueno, me parece un poco raro.} & \de{Naja, ich finde ihn etwas seltsam.} \\[3pt]
\textbf{B:} \es{¿Qué te parece su música?} & \de{Was hältst du von seiner Musik?} \\[3pt]
\textbf{A:} \es{Me parece bastante buena.} & \de{Ich finde sie ziemlich gut.} \\
\end{tabular}}
\end{merkbox}

\vspace{8pt}

\begin{multicols}{2}

% --- Zusammenfassung ---
\begin{grammatikbox}
\textbf{Zusammenfassung}

\vspace{3pt}
{\small
\textbf{1. parecer} (wie gustar):
\begin{itemize}[leftmargin=*, nosep, itemsep=0pt]
    \item Indir. Objektpronomen + parece(n)
    \item \es{Me parece bien.} = \de{Ich finde es gut.}
\end{itemize}

\vspace{3pt}
\textbf{2. conocer}:
\begin{itemize}[leftmargin=*, nosep, itemsep=0pt]
    \item yo = \es{conozco} (c → zc)
    \item + \es{a} bei Personen!
\end{itemize}}
\end{grammatikbox}

\columnbreak

% --- Typische Fehler ---
\begin{tippbox}
\footnotesize\textbf{Typische Fehler vermeiden:}

\vspace{2pt}
$\times$ \es{Conozco María.}\\
$\checkmark$ \es{Conozco \textbf{a} María.}

\vspace{2pt}
$\times$ \es{Yo conoco...}\\
$\checkmark$ \es{Yo \textbf{conozco}...}

\vspace{2pt}
$\times$ \es{Me parece simpáticos.}\\
$\checkmark$ \es{Me \textbf{parecen} simpáticos.}
\end{tippbox}

\end{multicols}

\vspace{8pt}

% ═══════════════════════════════════════════════════════════════
% ZUSATZAUFGABEN
% ═══════════════════════════════════════════════════════════════

\begin{zusatzbox}
{\small\textbf{Zusatzaufgaben für Schnelle}}

\vspace{4pt}

\begin{multicols}{2}
\footnotesize

\textbf{Z1: Interview}\\
Frage deinen Partner:
\begin{itemize}[leftmargin=*, nosep, itemsep=0pt]
    \item \es{¿Qué te parece nuestra clase?}
    \item \es{¿Conoces a alguien famoso?}
\end{itemize}

\vspace{4pt}

\textbf{Z2: Beschreibung}\\
Beschreibe eine Person mit \es{parecer}:\\
\es{Mi profesor/a me parece...}

\columnbreak

\textbf{Z3: Kettenspiel}\\
A: \es{¿Conoces a...?}\\
B: \es{Sí/No, (no) lo/la conozco.}\\
B fragt weiter...

\vspace{4pt}

\textbf{Z4: Kreativ}\\
Schreibe 3 Sätze: Was findest du gut/schlecht an deiner Schule?\\
\es{Me parece bien/mal que...}

\end{multicols}
\end{zusatzbox}

\end{document}
